\documentclass[11pt]{article}

%======================= Packages =======================
\usepackage[margin=1in]{geometry}
\usepackage{amsmath,amssymb,amsthm,mathtools}
\usepackage{bm}
\usepackage{enumitem}
\usepackage{hyperref}
\hypersetup{colorlinks=true,linkcolor=blue,citecolor=blue}

%======================= Math macros =======================
\newcommand{\E}{\mathbb{E}}
\newcommand{\Fset}{\mathcal{F}}
\newcommand{\R}{\mathbb{R}}
\newcommand{\Id}{\mathbb{1}}
\newcommand{\tp}{^{\prime}}
\DeclareMathOperator{\MSPE}{MSPE}
\DeclareMathOperator{\RMSE}{RMSE}
\DeclareMathOperator{\MAE}{MAE}
\DeclareMathOperator{\CRPS}{CRPS}
\DeclareMathOperator{\vecop}{vec}

\theoremstyle{remark}
\newtheorem{remark}{Remark}
\newtheorem*{note}{Note}

\title{Derivation Companion to\\
\emph{``Risk Scenarios and Macroeconomic Forecasts''}\\
\normalsize (Moran, Stevanovic \& Surprenant, BoC SWP 2025-28)}
\author{Equation-by-equation notes and derivations}
\date{\today}

\begin{document}
\maketitle

\begin{abstract}
\noindent This note reproduces every numbered equation of the paper, states the
assumptions under which it holds, and supplies the intermediate algebra that the
paper compresses. Numbering follows the paper: Eqs.~(1)--(5) build the
reduced-form VAR and the Waggoner--Zha (WZ) conditional forecast; Eq.~(6) and the
variance decomposition build the Baumeister--Kilian (BK) structural forecast; the
appendix equations (7)--(11) define the out-of-sample evaluation metrics. Three
places where the published display appears to contain a typo are flagged in
remarks.
\end{abstract}

\tableofcontents
\bigskip

%=========================================================
\section{Notation and conventions}
%=========================================================
Throughout, $y_t\in\R^{K}$ is the vector of $K$ endogenous variables, $p$ is the
lag order, and $H$ the maximum forecast horizon. The information set at time $t$
is $\Fset_t=\{y_s\}_{s=1}^{t}$. Innovations are (reduced-form) white noise:
\[
\E[u_t]=0,\qquad \E[u_tu_t\tp]=\Sigma_u,\qquad \E[u_tu_s\tp]=0\ \ (t\neq s).
\]
The companion objects are $Y_t\in\R^{Kp}$, $\mu\in\R^{Kp}$, $U_t\in\R^{Kp}$, and
the companion matrix $A\in\R^{Kp\times Kp}$. The selection matrix
\[
J:=\begin{bmatrix} I_K & 0_K & \cdots & 0_K\end{bmatrix}\in\R^{K\times Kp}
\]
extracts the top block, so that the following identities hold and are used
repeatedly:
\begin{equation}
y_t = J Y_t,\qquad \mu = J\tp\nu,\qquad U_t = J\tp u_t. \tag{$\star$}\label{eq:star}
\end{equation}
The reduced-form moving-average (impulse-response) coefficients are
$\Phi_j:=JA^jJ\tp$, and — once a structural impact matrix $D$ is identified — the
structural impulse responses are $\Theta_j:=\Phi_j D=JA^jJ\tp D$.

%=========================================================
\section{The reduced-form VAR and its companion form (Eqs.~1--2)}
%=========================================================
\subsection*{Eq.~(1): reduced form}
\begin{equation}
y_t = \nu + A_1 y_{t-1} + A_2 y_{t-2} + \cdots + A_p y_{t-p} + u_t. \tag{1}
\end{equation}
This is a $K$-dimensional linear autoregression of order $p$. ``Reduced form''
means $u_t$ collects statistical innovations with \emph{no} economic
interpretation; the $A_i$ summarize both own- and cross-dynamics. Estimation is
equation-by-equation OLS, which is efficient here because every equation has the
same right-hand-side regressors (the seemingly-unrelated-regression system
collapses to OLS).

\subsection*{Eq.~(2): companion (state-space) form}
\begin{equation}
Y_t = \mu + A Y_{t-1} + U_t. \tag{2}
\end{equation}

\paragraph{Derivation.} Stack the current and $p-1$ most recent vectors,
\[
Y_t:=\begin{bmatrix} y_t\\ y_{t-1}\\ \vdots\\ y_{t-p+1}\end{bmatrix},\qquad
\mu:=\begin{bmatrix} \nu\\ 0\\ \vdots\\ 0\end{bmatrix},\qquad
U_t:=\begin{bmatrix} u_t\\ 0\\ \vdots\\ 0\end{bmatrix},
\]
\[
A:=\begin{bmatrix}
A_1 & A_2 & \cdots & A_{p-1} & A_p\\
I_K & 0_K & \cdots & 0_K & 0_K\\
0_K & I_K & \cdots & 0_K & 0_K\\
\vdots & & \ddots & & \vdots\\
0_K & 0_K & \cdots & I_K & 0_K
\end{bmatrix}.
\]
Reading Eq.~(2) block-row by block-row: the \emph{first} block row is
$y_t=\nu+A_1y_{t-1}+\cdots+A_py_{t-p}+u_t$, i.e.\ Eq.~(1); block row $i\ge 2$ is
the trivial identity $y_{t-i+1}=y_{t-i+1}$, which merely carries lags forward in
the state. Hence (2) is exactly (1) written as a first-order system. The
eigenvalues of $A$ govern stability: covariance-stationarity requires all
eigenvalues to lie inside the unit circle.

%=========================================================
\section{Iterating forward: the $h$-step recursion (Eq.~3)}
%=========================================================
\begin{equation}
Y_{t+h}=\Bigg(\sum_{j=0}^{h-1}A^{j}\Bigg)\mu + A^{h}Y_t
+ \sum_{j=0}^{h-1}A^{j}U_{t+h-j}. \tag{3}
\end{equation}

\paragraph{Derivation (induction on $h$).}
\emph{Base case} $h=1$ is Eq.~(2) shifted forward:
$Y_{t+1}=\mu+AY_t+U_{t+1}$, which matches (3) since
$\sum_{j=0}^{0}A^{j}\mu=\mu$, $A^{1}Y_t=AY_t$, and
$\sum_{j=0}^{0}A^{j}U_{t+1-j}=U_{t+1}$.

\emph{Inductive step.} Assume (3) holds at horizon $h$. Apply (2) once more,
$Y_{t+h+1}=\mu+AY_{t+h}+U_{t+h+1}$, and substitute the inductive hypothesis:
\begin{align*}
Y_{t+h+1}
&=\mu+A\Bigg[\sum_{j=0}^{h-1}A^{j}\mu+A^{h}Y_t+\sum_{j=0}^{h-1}A^{j}U_{t+h-j}\Bigg]
+U_{t+h+1}\\[2pt]
&=\underbrace{\mu+\sum_{j=0}^{h-1}A^{j+1}\mu}_{=\sum_{j=0}^{h}A^{j}\mu}
+A^{h+1}Y_t
+\underbrace{\sum_{j=0}^{h-1}A^{j+1}U_{t+h-j}+U_{t+h+1}}_{=\sum_{j=0}^{h}A^{j}U_{t+h+1-j}}\\[2pt]
&=\sum_{j=0}^{h}A^{j}\mu+A^{h+1}Y_t+\sum_{j=0}^{h}A^{j}U_{t+h+1-j},
\end{align*}
which is (3) at horizon $h+1$. In the last shock sum, re-indexing
$j\mapsto j+1$ turns $\sum_{j=0}^{h-1}A^{j+1}U_{t+h-j}$ into
$\sum_{j=1}^{h}A^{j}U_{t+h+1-j}$, and the stray $U_{t+h+1}$ supplies the missing
$j=0$ term. 

%=========================================================
\section{Baseline forecast (Eq.~4)}
%=========================================================
\begin{equation}
\E\!\left[y_{t+h}\mid\Fset_t\right]
=\Bigg(\sum_{j=0}^{h-1}JA^{j}J\tp\Bigg)\nu + JA^{h}Y_t. \tag{4}
\end{equation}

\paragraph{Derivation.} Take $\E[\,\cdot\mid\Fset_t]$ of the companion recursion
(3). Every shock in the last term is dated $t+h-j\ge t+1$ for $j=0,\dots,h-1$,
hence is unforecastable: $\E[U_{t+h-j}\mid\Fset_t]=0$. Therefore
\[
\E[Y_{t+h}\mid\Fset_t]=\Bigg(\sum_{j=0}^{h-1}A^{j}\Bigg)\mu + A^{h}Y_t .
\]
Premultiply by $J$ to recover the top block $y_{t+h}=JY_{t+h}$, and use
$\mu=J\tp\nu$ from~\eqref{eq:star}:
\[
\E[y_{t+h}\mid\Fset_t]
=\Bigg(\sum_{j=0}^{h-1}JA^{j}\Bigg)J\tp\nu + JA^{h}Y_t
=\Bigg(\sum_{j=0}^{h-1}JA^{j}J\tp\Bigg)\nu + JA^{h}Y_t.
\]
Under quadratic (MSE) loss this conditional expectation is the optimal point
forecast; it is the \emph{baseline scenario}.

%=========================================================
\section{WZ conditional forecasts (Eq.~5 and the restriction $Ru=r$)}
%=========================================================
\subsection*{Eq.~(5): deviation of the path from the baseline}
\begin{equation}
y_{t+h}-\E\!\left[y_{t+h}\mid\Fset_t\right]
=\sum_{j=0}^{h-1}JA^{j}J\tp\,u_{t+h-j}
=\sum_{j=0}^{h-1}\Phi_j\,u_{t+h-j}. \tag{5}
\end{equation}

\paragraph{Derivation.} Subtract the conditional mean $\E[Y_{t+h}\mid\Fset_t]$
(the deterministic part of (3)) from (3) itself; only the shock term survives:
\[
Y_{t+h}-\E[Y_{t+h}\mid\Fset_t]=\sum_{j=0}^{h-1}A^{j}U_{t+h-j}.
\]
Premultiply by $J$ and use $U_{t+h-j}=J\tp u_{t+h-j}$ from~\eqref{eq:star}:
\[
y_{t+h}-\E[y_{t+h}\mid\Fset_t]
=\sum_{j=0}^{h-1}JA^{j}J\tp u_{t+h-j}
=\sum_{j=0}^{h-1}\Phi_j u_{t+h-j}. 
\]
\textbf{Economic reading.} Forcing a variable onto a scenario path (LHS
$\neq 0$) is \emph{equivalent} to imposing a particular realization of future
innovations (RHS). The scenario cannot happen ``for free''; it must be produced
by shocks, and those same shocks move every other variable through $\Phi_j$.

\subsection*{The stacked restriction $Ru=r$}
Stacking (5) over the constrained variables and horizons $h=1,\dots,H$ yields a
linear system in the future innovations. Let $M$ ($0<M<K$) be the number of
constrained variables. Following Blake--Mumtaz (2017),
\begin{equation}
Ru=r,\tag{5$'$}
\end{equation}
where:
\begin{itemize}[nosep]
\item $u:=\big[u_{1,t+1},\dots,u_{K,t+1},\ \dots,\ u_{1,t+H},\dots,u_{K,t+H}\big]\tp
\in\R^{KH}$ stacks all future innovations;
\item $r\in\R^{MH}$ is the LHS of (5) — the gap between the scenario values and
the baseline forecasts for the $M$ constrained variables over $H$ periods;
\item $R\in\R^{MH\times KH}$ collects the relevant rows of the MA coefficients
$\Phi_h$, i.e.\ the RHS map in (5).
\end{itemize}
Both $R$ and $r$ are known once the VAR is estimated.

\subsection*{The minimum-norm shock path (Doan--Litterman--Sims)}
Because $M<K$, system (5$'$) is \emph{underdetermined}: many shock sequences
reproduce the scenario. The paper selects the least-squares (minimum-norm)
solution.

\paragraph{Derivation (constrained minimization).} Solve
\[
\min_{u}\ \tfrac12\,u\tp u \quad\text{s.t.}\quad Ru=r.
\]
The Lagrangian is $\mathcal L=\tfrac12 u\tp u+\lambda\tp(r-Ru)$. The first-order
conditions are
\[
\partial_u\mathcal L=u-R\tp\lambda=0 \ \Rightarrow\ u=R\tp\lambda,
\qquad Ru=r.
\]
Substituting, $RR\tp\lambda=r$, so $\lambda=(RR\tp)^{-1}r$ (using that $R$ has
full row rank $MH$), and therefore
\begin{equation}
\hat u = R\tp(RR\tp)^{-1}r. \tag{5$''$}\label{eq:wzsol}
\end{equation}
Adding $\hat u$ back through (5) delivers the WZ conditional forecast for every
variable. 

\begin{remark}[dimensions of the inverse]
With $R$ of size $MH\times KH$ and $MH<KH$, the object $RR\tp$ is
$MH\times MH$ and invertible, whereas $R\tp R$ is $KH\times KH$ and
\emph{singular}. The paper's display writes $\hat u=R\tp(R\tp R)^{-1}r$; for the
stated dimensions this should read $(RR\tp)^{-1}$ as in~\eqref{eq:wzsol}. The two
orderings coincide only if one adopts the transposed convention for $R$ (i.e.\
storing $R$ as $KH\times MH$), in which case the roles of $RR\tp$ and $R\tp R$
swap. The substance is unchanged.
\end{remark}

\begin{remark}[``most likely'' shocks]
If innovations are standardized Gaussian, $u\sim N(0,I_{KH})$, then~\eqref{eq:wzsol}
is precisely the \emph{mode} (equivalently the mean) of the conditional
distribution of $u$ given $Ru=r$. Minimizing $u\tp u$ maximizes the Gaussian
density, which is why $\hat u$ is described as the most likely combination of
shocks consistent with the scenario.
\end{remark}

%=========================================================
\section{BK structural conditional forecasts (Eq.~6)}
%=========================================================
\subsection*{Eq.~(6): structural form}
\begin{equation}
y_t=\nu+A_1y_{t-1}+\cdots+A_py_{t-p}+D\epsilon_t. \tag{6}
\end{equation}
Here the reduced-form innovation is decomposed as $u_t=D\epsilon_t$, with
structural shocks that are white noise \emph{and mutually orthogonal},
normalized to unit variance:
\[
\E[\epsilon_t]=0,\qquad \E[\epsilon_t\epsilon_t\tp]=I_K
\ \Rightarrow\ \Sigma_u=\E[u_tu_t\tp]=D\,\E[\epsilon_t\epsilon_t\tp]\,D\tp=DD\tp.
\]

\subsection*{The BK forecast decomposition}
BK conditions on a path of \emph{structural shocks}
$\{\epsilon^{\text{scen}}_{t+h-j}\}$ rather than on variable values:
\begin{multline}
\E\!\left[y_{t+h}\,\Big|\,\{\epsilon_{t+h-j}=\epsilon^{\text{scen}}_{t+h-j}\}_{j=0}^{h-1},\Fset_t\right]
=\E\!\left[y_{t+h}\mid\Fset_t\right]
+\sum_{j=0}^{h-1}\Theta_j\,\epsilon^{\text{scen}}_{t+h-j}. \tag{6$'$}
\end{multline}

\paragraph{Derivation.} Start from the stochastic part of (3),
$\sum_{j=0}^{h-1}A^{j}U_{t+h-j}$, and project onto the top block as in (5):
\[
\sum_{j=0}^{h-1}A^{j}U_{t+h-j}
\ \xrightarrow{\ J\,\cdot\ }\
\sum_{j=0}^{h-1}JA^{j}J\tp u_{t+h-j}
=\sum_{j=0}^{h-1}JA^{j}J\tp D\,\epsilon_{t+h-j}
=\sum_{j=0}^{h-1}\Theta_j\,\epsilon_{t+h-j},
\]
using $u_t=D\epsilon_t$ and defining $\Theta_j:=JA^{j}J\tp D=\Phi_jD$ — the
matrix of structural impulse responses at horizon $j$. In the baseline these
future shocks average to zero; the scenario instead \emph{sets} them to
$\epsilon^{\text{scen}}$, and the induced deviation from the baseline is exactly
$\sum_{j}\Theta_j\epsilon^{\text{scen}}_{t+h-j}$, giving (6$'$).

\subsection*{Identification of $D$ via Cholesky}
\paragraph{Order condition (counting).} A symmetric $\Sigma_u$ has
$K(K+1)/2$ distinct entries, but $D$ has $K^2$ free entries. The identifying
equations $\Sigma_u=DD\tp$ therefore fall short by
\[
K^2-\tfrac{K(K+1)}{2}=\tfrac{K(K-1)}{2}
\]
restrictions. Imposing that $D$ be \emph{lower triangular} zeros out exactly its
$K(K-1)/2$ upper-triangular entries, which just closes the gap. With a positive
diagonal, the solution is the unique Cholesky factor $D=\mathrm{chol}(\hat\Sigma_u)$,
where $\hat\Sigma_u=\sum_{t}\hat u_t\hat u_t\tp/(T-Kp-1)$.

\paragraph{Timing.} A lower-triangular $D$ means that shock
$k$ affects variables $1,\dots,k-1$ only with a lag, but variables
$k,\dots,K$ contemporaneously. The recursive (causal) ordering of $y_t$ is
therefore material. The paper's ordering
\[
\text{oil}\to\text{US IP}\to\text{CPI}\to\text{housing}\to\text{unemp.}\to\text{disc.\ rate}\to\text{FX}
\]
places global variables first (small-open-economy block exogeneity), sticky
prices ahead of the activity/policy block, and the forward-looking exchange rate
last.

%=========================================================
\section{Forecast-error variance decomposition (Section 2.4)}
%=========================================================
\subsection*{The mean squared prediction error}
\begin{equation}
\MSPE(h)=\E\big[(y_{t+h}-\E y_{t+h})(y_{t+h}-\E y_{t+h})\tp\big]
=\sum_{j=0}^{h-1}\Theta_j\Theta_j\tp. \tag{7$'$}
\end{equation}

\paragraph{Derivation.} By (6$'$) the structural forecast error is
$y_{t+h}-\E y_{t+h}=\sum_{j=0}^{h-1}\Theta_j\epsilon_{t+h-j}$. Then
\[
\MSPE(h)=\E\!\left[\Big(\sum_{j=0}^{h-1}\Theta_j\epsilon_{t+h-j}\Big)
\Big(\sum_{i=0}^{h-1}\Theta_i\epsilon_{t+h-i}\Big)\tp\right]
=\sum_{j=0}^{h-1}\sum_{i=0}^{h-1}\Theta_j\,\E[\epsilon_{t+h-j}\epsilon_{t+h-i}\tp]\,\Theta_i\tp.
\]
Because the structural shocks are serially uncorrelated with identity
contemporaneous covariance, $\E[\epsilon_{t+h-j}\epsilon_{t+h-i}\tp]=I_K\,\delta_{ij}$
(Kronecker delta\footnote{The Kronecker delta $\delta_{ij}$ equals $1$ when $i=j$ and $0$ when $i\neq j$; multiplying by it keeps only the matched ($i=j$) terms and annihilates every cross term, collapsing the double sum to a single one.}), so all cross terms $i\neq j$ vanish and
$\MSPE(h)=\sum_{j=0}^{h-1}\Theta_j\Theta_j\tp$. 

\subsection*{Per-shock contributions and the relative decomposition}
Let $\theta_{k,l,j}$ be the $(k,l)$ entry of $\Theta_j$: the response of variable
$k$ to structural shock $l$ at horizon $j$. The $(k,k)$ diagonal of (7$'$) is
\[
\big[\MSPE(h)\big]_{kk}
=\sum_{j=0}^{h-1}\big[\Theta_j\Theta_j\tp\big]_{kk}
=\sum_{j=0}^{h-1}\sum_{l=1}^{K}\theta_{k,l,j}^{2}.
\]
Isolating a single shock $l$ gives its cumulative contribution to variable $k$,
\begin{equation}
\MSPE_{k,l}(h)=\theta_{k,l,0}^2+\theta_{k,l,1}^2+\cdots+\theta_{k,l,h-1}^2,
\qquad
\MSPE_{k}(h)=\sum_{l=1}^{K}\MSPE_{k,l}(h). \tag{7$''$}
\end{equation}
The reported \emph{relative} decomposition (Table~1) is the ratio
$\MSPE_{k,l}(h)/\MSPE_{k}(h)$ — the share of the $h$-step forecast-error variance
of variable $k$ attributable to shock $l$. Orthogonality of the $\epsilon$'s is
what makes these shares sum to one across $l$.

%=========================================================
\section{Appendix B: out-of-sample evaluation (Eqs.~7--11)}
%=========================================================
\subsection*{Eq.~(7): the AR benchmark}
\begin{equation}
y_{k,t}=\nu_k+a_1y_{k,t-1}+\cdots+a_py_{k,t-p}+u_{k,t}. \tag{7}
\end{equation}
This is (1) restricted to \emph{diagonal} coefficient matrices $A_1,\dots,A_p$:
each variable is forecast from its own lags only. It admits the same companion
form (2) and the same forecast formula (4), so it is a nested special case used
as the accuracy yardstick.

\subsection*{Eqs.~(8)--(10): point-forecast errors}
Define the trailing $h$-period average
$y^{(h)}_{k,t}:=\tfrac1h\sum_{j=1}^{h}y_{k,t-j+1}$. The $h$-step forecast error of
model $m$ is the realized-minus-predicted average,
\begin{equation}
e^{(h,m)}_{k,t}
=\frac1h\sum_{j=1}^{h}y^{(h)}_{k,t-j+1}
-\frac1h\sum_{j=1}^{h}\hat y^{(h,m)}_{k,t}, \tag{8}
\end{equation}
aggregated over the evaluation window $t=\underline t,\dots,\bar t$ into
\begin{align}
\RMSE(k,h,m)&=\sqrt{\frac{1}{\bar t-\underline t}\sum_{t=\underline t}^{\bar t}
\big(e^{(h,m)}_{k,t}\big)^2}, \tag{9}\\[2pt]
\MAE(k,h,m)&=\frac{1}{\bar t-\underline t}\sum_{t=\underline t}^{\bar t}
\big|e^{(h,m)}_{k,t}\big|. \tag{10}
\end{align}
The paper reports each of these \emph{relative} to the AR benchmark, so a value
$<1$ favors the VAR. RMSE penalizes large errors quadratically; MAE is robust to
outliers.

\subsection*{Eq.~(11): density forecasts via the CRPS}
Density forecasts are obtained by simulation: draw
$u_{t+h}\sim N(0,\hat\Sigma_u)$, iterate (1) forward $h=1,\dots,24$, and repeat
$5000$ times to build a predictive distribution at each horizon. These are scored
with the continuous ranked probability score. Let $\hat q^{(h,m)}_{\tau,k,t}$ be
the predictive $\tau$-quantile and let the quantile (``pinball'') loss be
\[
\rho_\tau(u):=u\big(\tau-\Id(u<0)\big).
\]
On a grid $0\le\tau_1<\cdots<\tau_N\le 1$,
\begin{equation}
\CRPS(k,h,m)=\frac{2}{N-1}\sum_{n=1}^{N}
\left(\frac{1}{\bar t-\underline t}\sum_{t=\underline t}^{\bar t}
\rho_{\tau_n}\!\big(y_{k,t}-\hat q^{(h,m)}_{\tau_n,k,t}\big)\right). \tag{11}
\end{equation}

\paragraph{Why this is a discretized CRPS.} For a predictive CDF $F$ and outcome
$y$, the CRPS admits the equivalent quantile representation
\[
\CRPS(F,y)=\int_{-\infty}^{\infty}\!\big(F(z)-\Id\{y\le z\}\big)^2\,dz
=2\int_0^1\rho_\tau\!\big(y-F^{-1}(\tau)\big)\,d\tau .
\]
Equation (11) replaces the integral $2\int_0^1(\cdot)\,d\tau$ by the trapezoidal
average $\tfrac{2}{N-1}\sum_n(\cdot)$ over the quantile grid, and the inner term
by the sample mean pinball loss over the evaluation window. This is the practical
appeal noted in the paper: only the predictive quantiles are needed, not the full
density. As with RMSE/MAE, ratios below one favor the VAR.

\begin{remark}[indicator in the pinball loss]
The proper pinball loss uses $\Id(u<0)$ (whether the realization falls below the
predicted quantile). The published display writes $\Id(u<\tau)$; the standard and
strictly-proper form is $\rho_\tau(u)=u(\tau-\Id(u<0))$, used above.
\end{remark}

%=========================================================
\section*{Summary of the logical chain}
\addcontentsline{toc}{section}{Summary of the logical chain}
\[
\underbrace{(1)\Rightarrow(2)}_{\text{companion}}
\ \Rightarrow\
\underbrace{(3)}_{\text{iterate}}
\ \Rightarrow\
\begin{cases}
\E[\cdot\mid\Fset_t]\ \Rightarrow\ (4)\ \text{baseline}\\[2pt]
\text{deviation}\ \Rightarrow\ (5)\ \Rightarrow\ Ru=r\ \Rightarrow\ \hat u=R\tp(RR\tp)^{-1}r\ \text{(WZ)}\\[2pt]
u=D\epsilon\ \Rightarrow\ (6),(6')\ \text{with}\ \Theta_j=\Phi_jD\ \text{(BK)}\\[2pt]
\Theta_j\Theta_j\tp\ \Rightarrow\ \text{variance decomposition (Table 1)}
\end{cases}
\]
Every scenario result in the paper is a projection of one master object — the
recursion (3) — either onto imposed variable values (WZ) or imposed structural
shocks (BK).

\end{document}